\section*{Bab \ref{ch:b3:bacteriology} --- Bakteriologi: Pertumbuhan, Fisiologi, dan Genetika}

\begin{solution}{exo:b3:bacteriology:1}
\omterm{def:b3:bacteriology:envelope}{Gram-positif}: membran sitoplasma yang dibungkus dinding \omterm{def:b3:bacteriology:envelope}{peptidoglikan} yang
tebal dan berlapis banyak serta disulami asam teikoat. \omterm{def:b3:bacteriology:envelope}{Gram-negatif}: sebuah
lapis \omterm{def:b3:bacteriology:envelope}{peptidoglikan} yang tipis di dalam sebuah periplasma, lalu sebuah \omterm{def:b3:bacteriology:envelope}{membran
luar} yang lembaran luarnya berupa \omterm{def:b3:bacteriology:envelope}{lipopolisakarida}, dengan porin. Kompleks ungu
kristal--iodinnya terperangkap di dalam dinding \omterm{def:b3:bacteriology:envelope}{Gram-positif} yang tebal dan
dikeringkan alkohol (ungu) tetapi tercuci keluar dari \omterm{def:b3:bacteriology:envelope}{Gram-negatif} lewat
dindingnya yang tipis dan \omterm{def:b3:bacteriology:envelope}{membran luarnya} yang terganggu, yang lalu menerima
warna tandingannya yang merah muda.
\end{solution}

\begin{solution}{exo:b3:bacteriology:2}
Sebesar $10^{6}$ kali lipat $= 2^{19.9}$: kira-kira $20$ generasi; $T = \qty{8}{h}/20
= \qty{24}{min}$; $\mu = \ln 2/\qty{0.4}{h} = \qty{1.7}{h^{-1}}$.
\end{solution}

\begin{solution}{exo:b3:bacteriology:3}
\omterm{def:b3:bacteriology:hgt}{Transformasi} --- Griffith (1928) serta Avery, MacLeod, dan McCarty
(1944): pneumokokus \omterm{def:b3:cancer-biology:hallmarks}{ganas} yang dimatikan panas mentransformasi pneumokokus
tak \omterm{def:b3:cancer-biology:hallmarks}{ganas} yang hidup, dan agennya adalah DNA. \omterm{def:b3:bacteriology:hgt}{Transduksi} --- Zinder dan
Lederberg (1952): fag P22 mengusung gen \emph{Salmonella} dari satu galur ke
galur lain lewat sebuah saringan yang menahan selnya. \omterm{def:b3:bacteriology:hgt}{Konjugasi} --- Lederberg
dan Tatum (1946): dua galur \emph{E.~coli} yang auksotrof memberi
rekombinan prototrof hanya ketika dicampurkan, sehingga menuntut sentuhan sel
(tabung-U Davis).
\end{solution}

\begin{solution}{exo:b3:bacteriology:4}
Penisilin: transpeptidase yang menaut-silangkan \omterm{def:b3:bacteriology:envelope}{peptidoglikan}, yang tak
dimiliki sel hewan. Streptomisin: subunit ribosom yang kecil,
yang bentuk bakterinya berbeda dari bentuk eukariot. Siprofloksasin: girase
DNA, yakni sebuah topoisomerase bakteri yang tak ada pada manusia. Rifampisin:
subunit $\beta$ RNA polimerase bakteri, yang tak dimiliki
enzim eukariotnya.
\end{solution}

\begin{solution}{exo:b3:bacteriology:5}
$S^{*} = K_{S}D/(\mu_{\max} - D)$, $X^{*} = Y(S_{0} - S^{*})$. $D = 0.2$:
$S^{*} = \qty{0.0067}{g/L}$, $X^{*} = \qty{0.80}{g/L}$. $D = 0.6$:
$S^{*} = 0.06$, $X^{*} = 0.78$. $D = 0.79$: $S^{*} = 1.58$, $X^{*} = 0.17$.
Tersapu keluar ketika $D = \mu(S_{0}) = 0.8\times 2/2.02 = \qty{0.79}{h^{-1}}$.
\end{solution}

\begin{solution}{exo:b3:bacteriology:6}
$(N/V)_{\text{amb}} = kA_{\text{amb}}/p$. Organ cahayanya: $0.01\times
6\times 10^{12}/100 = 6\times 10^{8}$ tiap liter $= 6\times 10^{5}$ tiap
mL. Air lautnya: $10\times 6\times 10^{12}/100 = 6\times 10^{11}$ tiap liter
$= 6\times 10^{8}$ tiap mL --- yakni seribu kali lebih tinggi, dan tak pernah
tercapai di laut.
\end{solution}

\begin{solution}{exo:b3:bacteriology:7}
\qty{25}{mm}: peka, MIC-nya rendah. \qty{12}{mm}: kepekaannya berkurang,
antara --- boleh jadi gagal pada dosis biasa. \qty{0}{mm}: kebal, dengan
pertumbuhannya sampai ke cakramnya. Obatnya berdifusi dari cakramnya sehingga
kepekatannya turun terhadap jarak dengan cara yang dapat diulang;
pertumbuhannya berhenti ketika kepekatannya menyamai MIC-nya, sehingga
jari-jari zona dan MIC-nya terkait lewat kurva kalibrasi yang ditegakkan dengan
galur yang MIC-nya diketahui.
\end{solution}

\begin{solution}{exo:b3:bacteriology:8}
\omterm{def:b3:genetic-engineering:tools}{Plasmid} konjugatifnya mengusung sebuah integron yang larik kasetnya menahan
beberapa gen kekebalan, ditambah transposon yang mengusung lebih banyak lagi;
\omterm{def:b3:genetic-engineering:tools}{plasmidnya} menyeberang dalam satu perkawinan lalu mengekspresikan semuanya.
Semuanya mengelompok karena seleksi bagi salah satu obatnya mempertahankan
seluruh \omterm{def:b3:genetic-engineering:tools}{plasmidnya}, sehingga gen yang menumpang padanya bertahan (seleksi
bersama); integron menangkap kasetnya; transposon melompat ke \omterm{def:b3:genetic-engineering:tools}{plasmidnya}; dan
\omterm{def:b3:genetic-engineering:tools}{plasmidnya}, bukan gennya, yang menjadi satuan pemindahannya.
\end{solution}

\begin{solution}{exo:b3:bacteriology:9}
Obatnya sendirian: $10^{-7}\times 10^{9} = 100$ sel kebal diharapkan. Dua
obat yang saling bebas: $10^{-14}\times 10^{9} = 10^{-5}$. Persister bukan
mutan: sebagian kecil sel yang tidur bertahan dari obat apa pun lalu tumbuh
kembali menjadi populasi yang peka ketika obatnya hilang, sehingga obat kedua
tak membantu; hanya pengobatan yang cukup lama untuk menangkapnya saat ia
bangun, atau obat yang aktif pada sel yang tidur, yang membantu.
\end{solution}

\begin{solution}{exo:b3:bacteriology:10}
$X > X^{*}$: pemakaian $\mu X/Y$ melampaui pasokannya, sehingga $S$ turun di
bawah $S^{*}$; lalu $\mu(S) < D$ dan $\mathrm{d}X/\mathrm{d}t < 0$, sehingga
$X$ turun kembali; seiring turunnya $X$, $S$ naik lagi. Kedua simpangannya
terbetulkan: keadaan tunaknya mantap. Dua spesies: masing-masing menuntut
$\mu_{i}(S) = D$ agar bertahan, dan hal itu menetapkan $S_{i}^{*}$-nya sendiri;
haranya hanya dapat bernilai satu, sehingga paling banyak satu spesies yang
bertahan. Spesies yang $S^{*}(D)$-nya lebih rendah mendorong $S$ ke bawah apa
yang dituntut spesies yang lain, dan yang lain itu tersapu keluar; spesies mana
itu dapat berubah dengan $D$ bila kurva Monodnya bersilangan.
\end{solution}

\begin{solution}{exo:b3:bacteriology:11}
Dengan kerapatan $\rho = N/V$, keadaan padamnya $A = p_{0}\rho/k$ ada selama
ia tetap di bawah $A_{\text{amb}}$, yakni $\rho < kA_{\text{amb}}/p_{0}$;
keadaan menyalanya $A = p_{1}\rho/k$ ada selama ia tetap di atasnya, yakni $\rho
> kA_{\text{amb}}/p_{1}$. Bagi $kA_{\text{amb}}/p_{1} < \rho <
kA_{\text{amb}}/p_{0}$ keduanya taat asas: sebuah populasi yang telah
menyala menjaga dirinya tetap menyala dengan keluarannya yang lebih tinggi,
sedangkan yang belum tetap padam. Kegunaannya: komitmen --- begitu sebuah
koloni memutuskan membangun biofilm atau melepaskan toksin, sebuah penurunan
kerapatan tak membatalkan keputusannya, dan sakelarnya tajam serta serentak,
bukan berjenjang.
\end{solution}

\begin{solution}{exo:b3:bacteriology:12}
(a) Dosis tinggi bagi kursus pendek: kadarnya tetap di atas MIC mutannya lalu
membunuh tipe liar maupun mutan langkah pertamanya --- dipihak, dengan batas
berupa ketoksikannya. (b) Dosis rendah bagi kursus panjang: kadarnya duduk di
dalam jendelanya berhari-hari lalu membiakkan kekebalan --- ditolak. (c)
Berhenti ketika gejalanya reda: penyintasnya adalah sel yang paling tak peka
dan kadarnya yang turun melewati jendelanya --- ditolak. (d) Gabungan:
sebuah sel harus kebal terhadap keduanya, yang tak dimuat semiliar sel
--- didukung.
\end{solution}

\begin{solution}{pb:b3:bacteriology:1}
\textbf{1.} $18$ generasi: $2^{18} = 2.6\times 10^{5}$ sel,
$2.6\times 10^{3}$ tiap mL.
\textbf{2.} $2\times 10^{9}\times 100 = 2\times 10^{11}$ sel $=
2^{37.5}$: sesudah \qty{12.5}{h}; massanya \qty{0.2}{g}.
\textbf{3.} $6\times 10^{27}/10^{-12} = 6\times 10^{39} = 2^{132}$ sel:
$132$ generasi, \qty{44}{h}. Hara, oksigen, dan limbahnya menghentikannya
dalam sehari; pertumbuhan eksponensial selalu singkat.
\textbf{4.} Sel stasioner telah membongkar ribosomnya dan menekan
enzim metabolismenya; semuanya harus dibangun kembali sebelum membelah.
Sel eksponensial datang dengan perlengkapan yang lengkap.
\textbf{5.} $\mu = \ln 2/\qty{0.333}{h} = \qty{2.08}{h^{-1}}$; pada $T =
\qty{60}{min}$, nilainya \qty{0.69}{h^{-1}}; sesudah \qty{6}{h} jumlahnya
$2^{6} = 64$ alih-alih $2^{18}$: yakni faktor $2^{12} \approx 4000$ lebih
sedikit.
\textbf{6.} $200/(0.1\times 10^{-6}) = 2\times 10^{9}$ sel hidup tiap
mL. Mikroskopnya juga menghitung sel yang mati dan sel yang tak akan tumbuh
pada pelatnya, dan sebuah gumpalan memberi satu koloni.
\textbf{7.} $D = 0.5/1 = \qty{0.5}{h^{-1}}$; $T = \ln 2/0.5 =
\qty{1.39}{h} = \qty{83}{min}$.
\textbf{8.} $S^{*} = 0.01\times 0.5/0.5 = \qty{0.01}{g/L}$; $X^{*} =
0.5\times 1.99 \approx \qty{1}{g/L} = 10^{12}$ sel tiap liter, $10^{9}$
tiap mL.
\textbf{9.} $0.5\times 10^{12} = 5\times 10^{11}$ sel tiap jam; $720/1.39
\approx 520$ generasi tiap bulan.
\textbf{10.} $D = 0.95$: $S^{*} = 0.01\times 0.95/0.05 = \qty{0.19}{g/L}$,
$X^{*} = 0.5\times 1.81 = \qty{0.9}{g/L}$. Pada $F = \qty{1.1}{L/h}$, $D
> \mu_{\max}$: tersapu keluar.
\textbf{11.} Mutannya menahan $S^{*} = 0.01\times 0.5/0.7 = \qty{0.0071}{g/L}$;
pada glukosa sebesar itu, galur aslinya tumbuh pada $1.0\times 0.0071/0.0171 =
\qty{0.42}{h^{-1}} < D$ lalu menurun sebagai $e^{-0.08t}$: tersapu keluar
sepanjang beberapa minggu. Mutannya mengambil alih.
\textbf{12.} $10^{-9}\times 10^{12} = 1000$ mutasi menguntungkan tiap
generasi: penyesuaiannya sinambung, dengan sebuah klon unggul baru menyapu
tiap beberapa ratus generasi --- kemostat adalah sebuah mesin evolusi.
\textbf{13.} $kA_{\text{amb}}/p = 0.05\times 6\times 10^{12}/500 =
6\times 10^{8}$ tiap liter $= 6\times 10^{5}$ sel tiap mL.
\textbf{14.} $10^{10}/6\times 10^{5} \approx 1.7\times 10^{4}$ kali lipat.
\textbf{15.} $20\times 6\times 10^{12}/500 = 2.4\times 10^{11}$ tiap liter,
yakni $2.4\times 10^{8}$ tiap mL --- enam orde besar di atas
$10^{2}$ tiap mL di laut: gelap.
\textbf{16.} Sepuluh kali lipat berarti $3.3$ pelipatduaan, \qty{6.6}{h}. Pada
$10^{9}$ tiap mL populasinya masih $1700$ kali di atas ambangnya: cahayanya
tetap menyala.
\textbf{17.} Di laut tak seorang pun melihat cahayanya dan cahaya itu tak
menyuapi inang mana pun: sebuah mutan yang gelap menghemat \qty{20}{\%} lalu
mengalahkan yang bersinar. Di dalam organnya, cumi-cuminya memilih cahaya ---
ia mengeluarkan atau tak menyuapi sel gelap yang curang --- sehingga ongkosnya
membeli sebuah rumah.
\textbf{18.} Halangilah reseptornya atau musnahkanlah \omterm{def:b3:bacteriology:quorum}{autopengimbasnya}
(pemadaman kuorum): bakterinya hidup tetapi tak pernah menggelar toksinnya, dan
sistem kekebalannya membersihkannya. Karena bakteri yang peka dan yang
``kebal'' tumbuh sama saja --- tak ada yang dibunuh --- keunggulan terpilih
karena lolos dari obatnya kecil, dan kekebalannya semestinya menyebar lebih
lambat daripada terhadap obat yang membunuh.
\textbf{19.} A: $10^{-8}\times 10^{9} = 10$; B: $100$; keduanya: $10^{-15}\times
10^{9} = 10^{-6}$.
\textbf{20.} $P_{0}(\text{A}) = e^{-10} = 5\times 10^{-5}$;
$P_{0}(\text{keduanya}) = e^{-10^{-6}} = 0.999999$.
\textbf{21.} $16 \to 8$: satu waktu paruh, \qty{4}{h}; $16 \to 1$: empat waktu
paruh, \qty{16}{h}; berada di dalam jendelanya dari \qty{4}{h} sampai
\qty{16}{h}, yakni \qty{12}{h} tiap dosisnya.
\textbf{22.} Dosis tiap \qty{12}{h}: berada di dalam jendelanya dari jam $4$
sampai jam $12$, yakni dua pertiga waktunya; sepanjang sepuluh hari, mutan
langkah pertama A --- sepuluh sel pada awalnya --- terpilih dan A sendirian
gagal.
\textbf{23.} Di atas \qty{8}{mg/L} sepanjang waktu: berilah dosis tiap satu
waktu paruh (\qty{4}{h}), atau sedosis yang memuncak pada \qty{64}{mg/L} tiap
\qty{12}{h} ($64 \to 8$ dalam tiga waktu paruh), atau sebuah infus sinambung.
Ongkosnya: ketoksikan puncak yang tinggi, atau enam dosis sehari yang tak akan
dipatuhi pasiennya.
\textbf{24.} $10^{-4}\times 10^{9} = 10^{5}$ persister bertahan apa pun
obatnya; ketika pengobatannya berhenti, semuanya bangun lalu menumbuhkan
kembali populasi yang sepenuhnya peka --- yakni sebuah kekambuhan. Pengobatan
harus berlangsung cukup lama untuk menangkapnya saat ia meninggalkan
ketiduran, dan itulah sebabnya tuberkulosis memakan enam bulan.
\textbf{25.} $10^{9}$ sel tiap mL di dalam kemostatnya; cahayanya menyala pada
$6\times 10^{5}$ sel tiap mL; $P_{0} = 0.999999$ bahwa kursus dua obatnya
tak berjumpa dengan sel kebal.
\end{solution}
